\chapter{NLPステートの再定義}
\epigraph{``The human nervous system cannot access two incompatible states at the same time. You cannot laugh hysterically and be terrified in the very same moment.''\\
（人間の神経系は、互いに相容れない二つのステートに同時にアクセスすることはできない。猛烈に爆笑しながら、同時に恐怖に震えることなど不可能なのだ。）}{Dr.~Richard Bandler}
人間のステートはDMN優位とCEN優位の二者択一にとどまらず、睡眠時など多様な状態を取りうる。ここでは、人間が取りうるステートを表す包括的な記号として $S$ を導入する。

ここではNLPで必要な日常的なDMN優位状態、通常のCEN優位状態、そしてCEN優位でありながら言語野のインタープリター（解釈器）が停止している基準状態（ZEROステート）を、それぞれ次のように定義する。
\begin{align}
    \State(\mathrm{DMN})  &\equiv \text{DMN優位（自己参照ナラティブ駆動）} \\
    \State(\mathrm{CEN})  &\equiv \text{CEN優位（通常の外部タスク駆動）} \\
    \State(\mathrm{ZERO}) &\equiv \text{CEN優位かつ言語野インタープリター停止}
\end{align}

また、これらのステートにおける現象を$::$で定義することとする。つまり、$\State(\mathrm{DMN})$のステートで、内的表象$Ve$が生じた時に確信$\Omega$に収束した時、

\begin{align}
    \State(\mathrm{DMN}) :: Ve \nbuc{}{} \Omega
\end{align}

と表す。

\subsection{変性意識状態（アルタードステート）の再定義：保留される確信 $\Omega$ と全神経系の探索ループ}

催眠やトランスの本質として長年議論されてきた「変性意識状態（Altered States of Consciousness: ASC / アルタードステート）」もまた、神経系疑似ベイズ推論機のダイナミクスによって明快に数学的再定義が可能である。

本稿では、アルタードステートを\textbf{「高次ナラティブ確信 $\Omega$ が確定せず、保留され続けている物理状態」}と定義する。

ミクロな感覚レベルにおいて解が確定しない空転状態（$e_{\mathrm{noise}} \nbscanning{\hat{\omega}} $）と同様に、マクロなナラティブの階層において、エビデンスのシークエンス $\boldsymbol{e}_{\mathrm{seq}}$ に対していかなる単一の確信 $\Omega$ にも収束できず、推論のゲートが開いたまま保留されている状態を次のように記述する。

\begin{align}
    &\boldsymbol{e}_{\mathrm{seq}} \nbscanning{\hat{\Omega}} \\
    &\equiv \State(ALTER)
    \label{eq:asc_definition}
\end{align}

この保留状態が引き起こす生物学的リアクションこそが、変性意識特有の精神・神経現象の核心である。

推論機が解を確定できないとき、脳は計算を放棄するのではなく、むしろ逆の挙動を示す。
確信 $\Omega$ という「確定した枠組み（解釈のゲシュタルト）」が成立しないため、神経系疑似ベイズ推論機は解を見出そうと、\textbf{エピソード記憶、身体感覚、無意識の連合ネットワークを含むらゆる神経系の全領域を総動員し、広範かつ並列的な仮説探索を激しく駆動させ続ける}のである。

とりわけ強力なアルタードステート（深いトランス状態）が誘発されるのは、\textbf{エゴにとって強烈な利益（profit）または壊滅的な脅威（threat）を予感させる極限の文脈において、既存の認知フレームではどうしても解 $\Omega$ が確定できない状況}に追い込まれた瞬間である。

エゴにとって利害が切迫しているため、神経系は計算を打ち切ることが許されない。
しかし、提示される言語パターン（ダブルバインド、逆説、文脈の急激な破壊など）が既存のどの確信 $\Omega_{\mathrm{ego}}$ にも適合しないとき、通常の思考ルーチンは完全にクラッシュする。

このとき生体には、以下のような特徴的兆候が現れる。
\begin{itemize}
  \item \textbf{DMNによる固定ナラティブの停止}：\\
  自己防衛の既成ストーリー（$A_{d,\mathrm{ego}}$）が解を見出せずフリーズし、自己と他者、内面と外界の境界認識が一時的に融解する。
  \item \textbf{被暗示性（Suggestibility）の極大化}：\\
  脳があらゆる回路を開いて猛烈に「解（確信）」を欲している飢餓状態にあるため、外部から投入される微小な言語シグナルや象徴的介入が、抵抗（エゴの検閲）を受けることなく、新たな確信 $\Omega$ として直接インストールされる。
\end{itemize}

リチャード・バンドラーが成田空港のTIMEX時計の話から「高級時計はすぐに狂っちまう」という一言を放った瞬間、私の脳内で起きたことの正体はまさにこれであった。
「高名なマスターの教え」という極めて重要な文脈（高サリエンス）において、時計への愛着や価値観というエゴの文脈が突如として不可解な矛盾に衝突させられ、解の確定が強制保留された（$\boldsymbol{e}_{\mathrm{seq}} \xrightarrow[]{\hat{\Omega}} \mid$）。
その瞬間、私の神経系は全記憶をスキャンする探索ループに叩き込まれ、無防備なトランス状態――すなわち変性意識状態へと転落したのである。


